\section{Implementare și Algoritmi}
\label{sec:algorithms}

Pentru a valida ipotezele teoretice, am implementat și comparat două arhitecturi distincte: un model serial de referință și algoritmul distribuit propus.

\subsection{Algoritmul Serial de Referință}
Ca punct de referință, am modelat abordarea centralizată standard. Un server central acționează ca un coordonator care interoghează secvențial sau concurent agenții de monitorizare.

Fluxul operațional este structurat în cicluri de monitorizare:
\begin{enumerate}
    \item \textbf{Colectarea (Polling):} Serverul recepționează valorile $v_i$ de la toate cele $N$ noduri active.
    \item \textbf{Calculul:} Serverul calculează suma totală $S_{total} = \sum v_i$ și derivă media aritmetică $\mu = S_{total} / N$.
\end{enumerate}

Deși acest algoritm oferă o precizie matematică absolută (în absența erorilor de rețea), el scalează liniar $\bigO{N}$ în ceea ce privește timpul de procesare și resursele necesare, devenind impracticabil în rețelele masive din cauza limitărilor de lățime de bandă la nivelul serverului \cite{van2017distributed}.

\subsection{Algoritmul Distribuit: Push-Pull Gossip}
Soluția propusă adoptă modelul \textit{SPMD (Single Program, Multiple Data)}, unde toate nodurile sunt echivalente funcțional. Arhitectura elimină coordonatorul central, bazându-se pe fire de execuție independente pentru emisia și recepția datelor.

Algoritmul utilizează protocolul \textbf{Push-Pull}, care este superior variantelor \textit{Push-Only} deoarece reduce timpul de convergență prin actualizarea simultană a stării ambelor noduri participante la o interacțiune. Pseudocodul \ref{alg:push-pull} descrie logica implementată pe fiecare nod.

\begin{algorithm}[H]
    \caption{Push-Pull Dynamic Average Consensus}
    \label{alg:push-pull}
    \begin{algorithmic}[1]
        \State Inițializează lista de vecini conform topologiei
        \State \textbf{Start Paralel: Thread Activ, Thread Pasiv}

        \State \textbf{Thread 1: Activ}
        \State \quad Măsoară încărcarea CPU curentă ($L_{curr}$)
        \State \quad \textbf{Atomic:} Actualizează media cu diferența $(L_{curr} - L_{last})$
        \State \quad Selectează aleatoriu un vecin $V$
        \State \quad Trimite media locală către $V$ (PUSH)
        \State \quad Primește media vecinului $V$ (PULL)
        \State \quad \textbf{Atomic:} $Media \leftarrow (Media + Media_{vecin}) / 2$

        \State \textbf{Thread 2: Pasiv}
        \State \quad Așteaptă conexiune de la un inițiator
        \State \quad Primește valoarea $Val_{in}$
        \State \quad \quad \textbf{Atomic:} Salvează media curentă în $R$
        \State \quad \quad \textbf{Atomic:} $Media \leftarrow (Media + Val_{in}) / 2$
        \State \quad Trimite $R$ înapoi către inițiator
    \end{algorithmic}
\end{algorithm}

\vspace{0.5em}
\noindent \textbf{Notă:} Operațiile \textit{Atomic} sunt protejate de un \textbf{Mutex} pentru a evita conflictele. Astfel, ne asigurăm că firele de execuție nu scriu simultan în aceeași variabilă, fapt ce ar duce la pierderea datelor și la un calcul greșit al mediei.

\subsection{Mecanismul de Eșec: "Negative Mass Injection"}
Un aspect tehnic crucial observat în implementarea \textit{Dynamic Consensus} este comportamentul sistemului la ieșirea bruscă a unui nod din rețea. Algoritmul se bazează pe injecția diferențelor ($\Delta$). 

Dacă un nod $A$ a partajat deja o valoare mare (ex. 100\%) cu rețeaua, "masa" sa a fost transferată parțial către vecini. Dacă nodul $A$ părăsește rețeaua (crash), el ia cu el valoarea sa curentă. Totuși, din perspectiva algoritmului distribuit, dacă nodul $A$ ar reveni cu valoarea 0, el ar injecta un $\Delta = 0 - 100 = -100$. 

În testele noastre de robustețe, am observat că dispariția nodurilor fără un protocol de \textit{graceful shutdown} poate duce la pierderea definitivă a unei cantități din suma globală ("Mass Loss"). Acest fenomen explică de ce, în anumite scenarii de avarie, media estimată a rețelei scade brusc, sistemul stabilizându-se la o nouă valoare de consens care reflectă "masa" rămasă în nodurile active \cite{kempe2003gossip}.